\section{Introduction}
Public procurement is a principal channel through which governments acquire artificial intelligence (AI) systems and services, and the market for such contracts is of direct interest to suppliers, buyers, and regulators \cite{hickok2022}. Quantifying this market requires deciding which procurement notices concern AI. In practice, this decision is usually operationalised with English keyword lists, as in studies that measure AI demand from online vacancies \cite{acemoglu2022,babina2024} and in public trackers of AI contracts. Keyword measures are transparent and cheap, but their accuracy is rarely reported against a reference annotation.

Two features of procurement data make this omission consequential. First, European notices published in Tenders Electronic Daily (TED) are written in 24 languages \cite{ahmia2018}; in the 54,160 IT and software award notices analysed here, only 6.1\% were published in English. An English dictionary may therefore miss AI purchases in a language-dependent way, which distorts cross-country comparisons. Second, an AI indicator produced by a dictionary or a classifier is often used as a regressor in a second-stage analysis. Recent work shows that plugging machine-generated variables into regressions yields biased estimates and invalid confidence intervals unless the measurement error is addressed \cite{battaglia2025,egami2023}, and that correction methods differ in finite samples \cite{audinet2025}.

This paper examines how the measurement of AI-relatedness in procurement notices affects three quantities: classification accuracy, the estimated share of AI notices, and the estimated association between AI-relatedness and single bidding, a standard indicator of weak competition in procurement \cite{fazekas2016}. We compare six scalable methods---two keyword dictionaries, a zero-shot natural language inference (NLI) classifier, an embedding-prototype score, and two supervised classifiers---against a codebook-based reference annotation of a stratified probability sample, and we apply design-based supervised learning (DSL) \cite{egami2023} to obtain prevalence and regression estimates that account for classifier error. All hypotheses, the sample design, and the codebook were fixed before annotation in a pre-specified analysis plan.

The contributions are threefold. (i) We provide, to our knowledge, the first language-stratified accuracy assessment of AI identification in multilingual procurement notices. (ii) We show how prevalence and competition estimates change with the measurement method and after DSL correction. (iii) We release code, prompts, and notice-level derived data so that the procedure can be reused for procurement monitoring. The contribution is methodological and empirical; no new learning algorithm is proposed.

\section{Related Work}
\subsubsection{Text-based measurement of AI demand.} Labour-economics studies identify AI demand through AI-related terms in job postings \cite{acemoglu2022,babina2024}, and large language models (LLMs) have been used to score occupational exposure to AI \cite{eloundou2024}. Text-as-data surveys emphasise that dictionary and supervised measures trade transparency against recall and that their error should be quantified \cite{gentzkow2019}. These studies mainly analyse English text and seldom report precision and recall against an independent annotation.

\subsubsection{Procurement analytics.} TED data support corpus construction \cite{ahmia2018}, national award databases \cite{potin2023}, process mining \cite{nai2023}, and CPV code assignment \cite{siciliani2023}. Single bidding and related indicators are widely used to assess competition and integrity risk \cite{fazekas2016,decarolis2022}, while the number of bidders has proven difficult to predict from notice features \cite{lozamencia2013}. To our knowledge, no study has evaluated how AI contracts are identified in multilingual TED notices.

\subsubsection{LLM annotation and inference with predicted variables.} LLMs can match or exceed crowd workers on several annotation tasks \cite{gilardi2023,tornberg2023}, but performance varies widely across tasks \cite{ziems2024}. Prediction-powered inference (PPI) \cite{angelopoulos2023} and DSL \cite{egami2023} combine a small labelled sample with abundant predictions to obtain valid estimates; benchmark evidence suggests that DSL often reduces bias more than PPI, although not uniformly \cite{audinet2025}. Multilingual sentence encoders \cite{wang2024me5,reimers2020} and NLI-based zero-shot classification \cite{yin2019,he2023} offer language-independent alternatives to dictionaries, with performance that varies by task and language \cite{muennighoff2023,conneau2020}.
